\section*{Experiencia Técnica}

\textbf{Pipelines de Inferencia Bayesiana}
\begin{itemize}
    \item Diseñé y desarrollé \href{https://github.com/seap-udea/PRisma}{PRisma}, un framework modular de inferencia bayesiana para investigación en exoplanetas, integrando modelos forward, funciones de verosimilitud basadas en KDE e inferencia mediante nested sampling (\texttt{dynesty}). Implementé ejecución paralela en clústeres HPC, flujos de trabajo reproducibles y campañas de inferencia versionadas, permitiendo análisis a gran escala y respaldando investigación como primer autor.
\end{itemize}
\textit{Python, Inferencia bayesiana, Nested sampling (dynesty), Computación paralela, HPC, Cómputo científico, Investigación reproducible}

\vspace{6pt}
\textbf{Minería de Datos y Machine Learning Aplicados}
\begin{itemize}
    \item Apliqué flujos completos de ciencia de datos sobre conjuntos de datos astronómicos, incluyendo consultas SQL/ADQL, preprocesamiento, visualización, ingeniería de características y aprendizaje automático supervisado y no supervisado. Desarrollé modelos predictivos y de agrupamiento utilizando el ecosistema científico de Python en proyectos orientados a investigación.
\end{itemize}
\textit{SQL, ADQL, scikit-learn, Análisis de datos, Clasificación, Clustering, Machine learning, Visualización de datos}

\vspace{6pt}
\textbf{Ingeniería de Software Científico}
\begin{itemize}
    \item Co-desarrollé y mantuve \href{https://github.com/seap-udea/pryngles}{Pryngles}, un paquete open source en Python para simulaciones fotométricas de sistemas planetarios. Diseñé y evolucioné su API pública, implementé infraestructura de pruebas y documentación, y apliqué prácticas modernas de ingeniería de software, incluyendo arquitectura modular, flujos colaborativos con Git, desarrollo reproducible y mantenibilidad a largo plazo.
\end{itemize}
\textit{Python, Arquitectura de software, Diseño de API, Desarrollo open-source, Git / GitHub, Pruebas unitarias, Documentación, Investigación reproducible}
